%==============================================================================
% Section 8: Conclusion
%==============================================================================
\section{Conclusion}
\label{sec:conclusion}

We have reorganised the homogeneous $\Sthree\times\Sone$ minisuperspace geometry
of Einstein--Cartan + Nieh--Yan theory along the final irreducible basis.  As a
result, the following five structures close as a single theoretical framework:
\begin{itemize}
  \item Pontryagin protection in the AX/VT torsion modes,
  \item completion of the spin-2 quintet by the full shear tensor,
  \item physical/auxiliary decomposition of the twist-mixing system,
  \item quadratic decoupling in the final basis, and
  \item a unified geometric Landau-EFT in the invariant basis.
\end{itemize}

In particular, the central message of this paper can be distilled into one
sentence:
\begin{quote}
  \itshape
  The homogeneous geometric sector of the minisuperspace is a completed
  geometric Landau-EFT, organised along the spin-0/1/2 irreducible
  representations.
\end{quote}

From this conclusion, the present model establishes one closed framework for the
classical homogeneous geometric sector.  On the other hand, the next challenges
--- thermal interpretation, local degrees of freedom, and interface analysis ---
are left to future work.  This paper provides the theoretical junction that
delineates those branches.
